Les anions thiosulfate $\ce{S2O^{2-}_3}$ sont capables de réduire l'iode
$\ce{I2}$ en ion iodure $\ce{I-}$.On utilise cette réaction pour doser
$\ce{I2}$.

On introduit dans un bécher un volume $V=\qty{20}{\mL}$ d'une solution de l'iode
$\ce{I2}$ de concentration $C$, et par une burette on ajoute progressivement une
solution de thiosulfate de sodium $(\ce{2Na^{+}}, \ce{S2O^{2-}_3})$ de
concentration $C_1=\qty{8.0e-2}{\mol\per\L}$, on observe un changement de
couleur de solution lorsqu'on ajoute un volume $V_{1E}=\qty{15.8}{\mL}$.

\begin{donnees}
    \begin{itemize}
        \item $\ce{S4O^{2-}_{6(aq)}}\ttslash{}\ce{S2O^{2-}_{3(aq)}}$ et
              $\ce{I_{2(aq)}}\ttslash{}\ce{I^-_{(aq)}}$
    \end{itemize}
\end{donnees}
\begin{enumerate}
    \item Quel est le but d'un dosage.
    \item Faire un schéma annoté du dispositif expérimental.
    \item Déterminer le réactif titrant et le réactif titré dans ce dosage
    \item Écrire les deux demi-équations électroniques relatives aux couples mis
          en jeu.
    \item Déduire l'équation de la réaction de dosage.
    \item Écrire le tableau d'avancement de cette réaction
    \item Quel est le type de ce dosage, comment repérer expérimentalement
          l'équivalence~?
    \item Calculer la concentration $C$ de l'iode $\ce{I2}$.
    \item Faire l'inventaire des ions présents dans le mélange à l'équivalence
          et calculer leurs concentrations.
    \item Donner l'expression de la conductivité $\sigma$ à l'équivalence.
\end{enumerate}
